% 04-profiler.tex
\section{剖析器子系统：OpRecord、GpuTimerPair 与后端栈}
\label{sec:profiler}

\subsection{目标和安置}

分析器必须满足三个相互冲突的目标：(i) 禁用时接近零开销，(ii) 启用时准确的 CPU+GPU 计时，以及 (iii) 后端的可插拔性（CPU 跟踪、CUDA CUPTI 活动、NVTX 范围、ITT 任务）。 TensorPlay 将核心置于 \texttt{p10} 中，因此 \texttt{tpx}（引擎）和内核都可以在不依赖 Python 的情况下进行检测。 Python 层 (\texttt{tensorplay.profiler}) 仅切换原子并耗尽缓冲区。放置遵循依赖规则：\texttt{Profiler.h} 是单个标头； \texttt{tpx} 和 \texttt{\_C} 包含它，而不是相反。

热路径是每个调度的操作加载一个 \texttt{std::atomic<bool>} 负载，当请求形状或站点捕获时，每个加载一个。这与每次调用 \texttt{GradMode} 和 \texttt{InferenceMode} 时已经发出的防护类相匹配。所有其他工作（形状复制、名称竞技场推送、GPU 事件记录）都位于该分支后面。

\subsection{核心数据模型：Event、GpuActivity 和 MemEvent}

定义了三个会话生命周期记录：

\begin{table}[H]
\centering
\caption{\texttt{Profiler.h} 中的探查器记录类型。}
\label{tab:profrecords}
\small
\begin{tabular}{lll}
\toprule
Type & 线路 & 目的 \\
\midrule
\texttt{Event} &  & 每个调度操作/用户范围/后向节点一个 \\
\texttt{MemEvent} &  & 当设置 \texttt{g\_mem\_capture} 时，每个分配器分配/释放一个 \\
\texttt{GpuActivity} &  & 每个 CUPTI 内核 / memcpy / memset / 运行时 / 驱动程序记录一个 \\
\bottomrule
\end{tabular}
\end{table}

\subsubsection{事件字段}

\begin{lstlisting}[caption={Event layout.}, label=lst:eventfields, style=cppstyle]
struct Event {
  const char* name;        // borrowed for kOp (static literal), arena-owned for kUser/kBackward
  uint64_t start_ns;       // CLOCK_MONOTONIC (steady_clock)
  uint64_t end_ns;
  uint64_t tid;            // hashed std::this_thread::get_id()
  EventKind kind;          // kOp='o', kUser='u', kBackward='b' (Profiler.h:43)
  std::shared_ptr<const ShapeVec> shapes; // record_shapes only
  std::shared_ptr<const DtypeVec> dtypes;
  uint32_t site_id = kNoSite;   // Python call site (with_stack)
  uint32_t stack_id = kNoSite;  // full frame chain (with_stack)
  void* gpu_start = nullptr;    // opaque cudaEvent_t pair (USE_CUDA)
  void* gpu_end = nullptr;
  float gpu_ms = -1.f;          // filled by gpu_resolve_all or summed kernel time in gpu_trace
  int32_t kernel_count = 0;     // correlated kernel count in gpu_trace mode
  int64_t out_bytes = -1;       // numel * itemsize for Tensor-returning ops
};
\end{lstlisting}

\texttt{EventKind} at 区分分派操作员 (\texttt{kOp})、通过 \texttt{record\_function} (\texttt{kUser}) 的用户注释以及 autograd 引擎的 \texttt{\_\_backward\_\_} 范围 (\texttt{kBackward})。 \texttt{site\_id} 和 \texttt{stack\_id} 通过 \texttt{site\_string} 和 \texttt{stack\_frames} 解决，并保持有效直到下一个 \texttt{profiler\_start}。

\texttt{GpuActivity} 添加 \texttt{correlation}、\texttt{external\_id}、\texttt{thread\_id}、\texttt{cbid}、\texttt{device}、\texttt{stream}、\texttt{kind}、\texttt{bytes}、\texttt{copy\_kind}、\texttt{value}。 \texttt{kNoExt = 0xffffffffffffffff} 标记不相关的活动。

\subsubsection{OpRecord 字段和生命周期}

\begin{lstlisting}[caption={OpRecord handle.}, label=lst:oprecorddecl, style=cppstyle]
struct TENSORPLAY_API OpRecord {
  explicit OpRecord(const char* name, EventKind kind = EventKind::kOp);
  explicit OpRecord(const std::string& name, EventKind kind = EventKind::kUser);
  ~OpRecord();
  void set_io_meta(ShapeVec&& shapes, DtypeVec&& dtypes);
  void set_output_bytes(int64_t nbytes);
  void set_shapes(ShapeVec&& shapes);
private:
  friend struct GpuTimerPair;
  void begin(const char* static_name, const std::string* owned_name, EventKind kind);
  uint64_t start_ns_ = 0;
  size_t slot_ = 0;
  bool live_ = false;
  bool counted_ = false;
  bool nvtx_open_ = false;
  bool itt_open_ = false;
  bool trace_pushed_ = false;
};
\end{lstlisting}

字段语义，按声明顺序：

\begin{table}[H]
\centering
\caption{OpRecord 私有状态。}
\label{tab:oprecordstate}
\small
\begin{tabular}{lp{9.2cm}}
\toprule
场地 & 意义 \\
\midrule
\texttt{start\_ns\_} & \texttt{steady\_clock} 在 \texttt{begin} 拍摄的快照 \\
\texttt{slot\_} & \texttt{g\_events} 向量的索引；稳定，因为向量是仅附加的 \\
\texttt{live\_} & 当且仅当此实例拥有有效插槽时为 true；守卫 \texttt{set\_io\_meta} 和 \texttt{\textasciitilde OpRecord} \\
\texttt{counted\_} & True iff，当且仅当 \texttt{g\_inflight} 增加；守卫 \texttt{release\_inflight} \\
\texttt{nvtx\_open\_} & 真当且仅当 \texttt{nvtx\_span\_begin} 被发射；即使 \texttt{live\_} 为 false，也会在析构函数中关闭 \\
\texttt{itt\_open\_} & \texttt{itt\_span\_begin} 相同 \\
\texttt{trace\_pushed\_} & True 当且仅当 \texttt{cupti\_push\_ext(slot\_)} 成功；出现在 \texttt{GpuTimerPair::close} \\
\bottomrule
\end{tabular}
\end{table}

生命周期，具体路径：线迹：

\begin{enumerate}[leftmargin=1.4em]
\item 构造调用\texttt{OpRecord::begin}。首先，它通过松弛负载检查 \texttt{g\_emit\_nvtx} 和 \texttt{g\_emit\_itt}（因此 NVTX/ITT 范围可以在没有 \texttt{kOp} 分析会话的情况下启动，如所述）。然后它检查 \texttt{g\_active};如果为 false，则立即返回 \texttt{live\_=false}：禁用成本是一个获取负载和一个预测的未采用分支。
\item 如果处于活动状态，它会执行 \texttt{g\_inflight.fetch\_add(1)}、快照 \texttt{now\_ns}（通过 \texttt{steady\_clock}），然后锁定 \texttt{g\_mutex}。 \texttt{g\_active} 的第二次宽松检查处理 \texttt{profiler\_stop} 清除两次检查之间的标志的竞争。
\item 它构建一个 \texttt{Event}： \texttt{name} 要么是静态文字（\texttt{kOp} 路径），要么是 \texttt{g\_name\_arena} 中的竞技场副本（对于 \texttt{kUser}/\texttt{kBackward}）。 \texttt{tid} 来自 \texttt{this\_thread\_id} （定义为 \texttt{std::this\_thread::get\_id} 的 \texttt{thread\_local} 哈希值）。如果设置了 \texttt{t\_pending\_site.valid}，则其 \texttt{stack\_id} 将被消耗并清除，因此内部复合操作不会继承最外层的 Python 站点。
\item 如果是 \texttt{nvtx\_on || itt\_on}，则调用 \texttt{nvtx\_span\_begin(e.name)} 和 \texttt{itt\_span\_begin(e.name)} 并设置 \texttt{itt\_open\_}。然后是 \texttt{g\_events->push\_back(e)} 和 \texttt{slot\_ = size-1}、\texttt{start\_ns\_ = start}、\texttt{live\_=true}。
\item 如果 \texttt{live\_} 和槽在边界内，则 \texttt{set\_io\_meta} 和 \texttt{set\_output\_bytes} 锁定 \texttt{g\_mutex} 并写入 \texttt{(*g\_events)[slot\_]}。 \texttt{set\_shapes} 是生成的调度程序的仅形状重载，其中数据类型保持未设置。
\item 销毁总是先关闭 NVTX/ITT，然后如果 \texttt{!live\_} 返回。否则，它记录 \texttt{end = now\_ns}，锁定 \texttt{g\_mutex}，如果槽仍然存在则写入 \texttt{(*g\_events)[slot\_].end\_ns = end}（如果 \texttt{profiler\_stop} 已经清除了向量，则静默删除），然后调用 \texttt{release\_inflight} 执行 \texttt{g\_inflight.fetch\_sub(1)} 并在计数达到零时通知 \texttt{g\_inflight\_cv}。
\end{enumerate}

\texttt{g\_inflight} 协议对于正确性至关重要：\texttt{profiler\_stop} 首先存储 \texttt{g\_active=false}，然后等待 \texttt{g\_inflight\_cv} 直到 \texttt{g\_inflight==0}，然后取出 \texttt{g\_mutex} 并将 \texttt{g\_events} 移出。这可以防止 CUDA 或 CPU 工作线程在 \texttt{gpu\_resolve\_all} 对向量进行快照后追加 \texttt{gpu\_ms} 或 \texttt{end\_ns}。

\begin{figure}[H]
\centering
\begin{tikzpicture}[node distance=0.7cm, font=\small]
\node[libbox, fill=blue!6, draw=tpblue, minimum width=13cm, minimum height=0.6cm] (cfg) {\texttt{ g\_capture\_sites} atomics \quad};
\node[libbox, fill=orange!7, draw=tporange, below=0.5cm of cfg, minimum width=6.2cm, minimum height=1.8cm, align=left] (begin) {\textbf{OpRecord::begin}\\ \texttt{now\_ns} + lock + push Event\\ slot\_ = events.size-1};
\node[libbox, fill=orange!7, draw=tporange, right=0.6cm of begin, minimum width=6.2cm, minimum height=1.8cm, align=left] (dtor) {\textbf{\textasciitilde OpRecord}\\ \texttt{now\_ns} + lock + end\_ns\\ \texttt{g\_inflight--}};
\node[libbox, fill=purple!7, draw=tppurple, below=0.6cm of begin, minimum width=6.2cm, minimum height=1.1cm] (shape) {\texttt{ set\_output\_bytes}};
\node[libbox, fill=green!7, draw=tpgreen, below=0.6cm of dtor, minimum width=6.2cm, minimum height=1.1cm] (gpu) {\texttt{ close}};
\node[libbox, fill=blue!6, draw=tpblue, below=0.6cm of shape, minimum width=13cm, minimum height=0.7cm] (arena) {\texttt{g\_events vector + g\_name\_arena + g\_site\_table} \quad guarded by \texttt{g\_mutex}};
\draw[arrow] (cfg) -- (begin);
\draw[arrow] (cfg) -- (dtor);
\draw[arrow] (begin) -- (shape);
\draw[arrow] (shape) -- (arena);
\draw[arrow] (dtor) -- (gpu);
\draw[arrow] (gpu) -- (arena);
\end{tikzpicture}
\caption{OpRecord 生命周期和存储。每个分派的操作都会创建一个 \texttt{OpRecord} 重分派漏斗； \texttt{g\_inflight} 阻止 \texttt{profiler\_stop} 与析构函数竞争。}
\label{fig:oprecordlife}
\end{figure}

\subsection{存储：仅附加向量作为逻辑 RingBuffer}

原始设计文档将此组件称为 \texttt{RingBuffer}（草案中为 \texttt{RingBuffer.h}）。交付的实现实现了与由互斥锁和稳定槽索引保护的仅附加向量相同的逻辑缓冲区。具体状态为：

\begin{lstlisting}[caption={session storage.}, label=lst:profstorage, style=cppstyle]
std::mutex g_mutex;
std::vector<Event>* g_events = nullptr;              // append-only, slots stable
std::vector<MemEvent>* g_mem_events = nullptr;
std::deque<std::string>* g_name_arena = nullptr;     // owns user/backward names
std::vector<std::pair<std::string,int>>* g_site_table = nullptr;
std::unordered_map<std::string,uint32_t>* g_site_index = nullptr;
std::deque<std::string>* g_frame_arena = nullptr;
std::unordered_map<std::string,uint32_t>* g_frame_index = nullptr;
std::vector<std::vector<uint32_t>>* g_stack_table = nullptr;
std::deque<ProfFrame>* g_frame_store = nullptr;
std::atomic<uint64_t> g_inflight{0};
std::mutex g_inflight_mutex;
std::condition_variable g_inflight_cv;
\end{lstlisting}

为什么是向量而不是固定环：在整个会话中，槽由 \texttt{OpRecord::slot\_} 和 \texttt{GpuTimerPair} 的外部相关 ID（\texttt{slot\_} 转换为 \texttt{uint64\_t} at）引用。覆盖的环将使这些索引无效。结合层通过 \texttt{std::move(*g\_events)} 排出 \texttt{profiler\_stop} 中的向量，然后通过 \texttt{start\_ns} 排出 \texttt{stable\_sort} 中的向量。 arena (\texttt{g\_name\_arena}) 一直保持活动状态，直到下一个 \texttt{profiler\_start}，因为用户跨度的 \texttt{Event::name} 是借用的指针（注释）。

待处理的 Python 站点处理是线程本地的：

\begin{lstlisting}[caption={per-thread pending site.}, label=lst:pendingsite, style=cppstyle]
thread_local struct {
  bool valid = false;
  uint32_t site_id = Event::kNoSite;
  uint32_t stack_id = Event::kNoSite;
} t_pending_site;
\end{lstlisting}

\texttt{set\_python\_site} 和 \texttt{set\_python\_stack} 在持有 GIL 时由绑定层调用；他们通过 \texttt{intern\_site} 和 \texttt{intern\_stack} 进行实习，并通过 \texttt{g\_site\_index} 和 \texttt{g\_frame\_index} 进行重复数据删除。该线程上的下一个 \texttt{OpRecord::begin} 消耗待处理站点并清除 \texttt{valid}，因此内部复合操作不记录任何站点，而不是重复最外层的调用。

用户注释在 (\texttt{thread\_local std::vector<size\_t> t\_user\_stack}) 处有自己的 LIFO 堆栈，推送到 \texttt{user\_span\_begin} 并弹出到 \texttt{user\_span\_end}。两者都在 \texttt{g\_mutex} 下操纵 \texttt{g\_events} 并在推/弹出周围发出 NVTX。

保留的逻辑 RingBuffer 属性：热路径上没有超出向量摊销增长的分配（op 事件为 64 字节加上可选的共享\_ptr 有效负载），并且缓冲区在会话停止时恰好耗尽一次。

\subsection{配置原子：ProfilerConfig Surface}

配置是命名空间 \texttt{tensorplay::prof} 中的一小组免费原子——每个功能开关一个：

\begin{table}[H]
\centering
\caption{探查器配置原子。草稿名称 $\rightarrow$ 发货符号。}
\label{tab:profatomics}
\small
\begin{tabular}{llll}
\toprule
草稿 \texttt{ProfilerConfig} & 运送原子 & 线路 & 防护功能 \\
\midrule
\texttt{enabled\_} & \texttt{g\_active} & , & 全部录音 \\
\texttt{with\_shapes\_} & \texttt{g\_capture\_shapes} & , & 重新调度中的 \texttt{set\_io\_meta} \\
\texttt{with\_stack\_} & \texttt{g\_capture\_sites} & , & \texttt{stack} 收养 \\
\texttt{capture\_memory\_} & \texttt{g\_mem\_capture} & , & \texttt{free} \\
\texttt{gpu\_timing\_} & \texttt{g\_gpu\_timing} & , & 每个 CUDA 操作 \texttt{cudaEvent} 对 \\
\texttt{gpu\_trace\_} & \texttt{g\_gpu\_trace} & , & CUPTI活动合集 \\
\texttt{emit\_nvtx\_} & \texttt{g\_emit\_nvtx} & , & 每个跨度 \texttt{Pop} \\
\texttt{emit\_itt\_} & \texttt{g\_emit\_itt} & , & 每个跨度 \texttt{end} \\
\bottomrule
\end{tabular}
\end{table}

所有都是 \texttt{std::atomic<bool>}，存储上有 \texttt{memory\_order\_release}（\texttt{Profiler.cpp}，via），热路径负载上有 \texttt{memory\_order\_acquire}。 \texttt{g\_mutex} 关键部分内出现两个额外的宽松检查，以处理开始/停止竞赛，而无需额外的同步。 Python 入口点将 \texttt{profile(record\_shapes,with\_stack,gpu\_timing,gpu\_trace,profile\_memory)} 存储为这些原子以及对 \texttt{profiler\_start}、\texttt{profiler\_start\_with\_shapes}、\texttt{profiler\_start\_full} 和 \texttt{cupti\_start} 的调用。

\texttt{with\_stack} 在 Python 层中隐含 \texttt{record\_shapes} ； C++ 层保持它们独立，但当设置 \texttt{with\_stack} 时，绑定的 \texttt{\_profiler\_start} 会分支到 \texttt{profiler\_start\_full}，这会将 \texttt{g\_capture\_shapes} 和 \texttt{g\_capture\_sites} 设置为 true。

\subsection{GpuTimerPair：cudaEvent 与\ CUPTI 选择}

将 \texttt{GpuTimerPair} 声明为池支持的 \texttt{cudaEvent\_t} 对加上 CUPTI 相关括号：

\begin{lstlisting}[caption={GpuTimerPair handle.}, label=lst:gputimerpair, style=cppstyle]
struct TENSORPLAY_API GpuTimerPair {
  explicit GpuTimerPair(OpRecord& rec) : rec_(rec) {}
  ~GpuTimerPair();
  void arm(const Device& device);  // records start on a CUDA op's stream
  void close();                    // records the end event
private:
  OpRecord& rec_;
  void* gpu_start_ = nullptr;
  void* gpu_stream_ = nullptr;
  int gpu_device_ = -1;
  bool armed_ = false;
};
\end{lstlisting}

该实现位于 \texttt{ProfilerGpu.cpp} 中。对于 \texttt{!USE\_CUDA} ，它在生成的代码链接处编译为无操作，没有 \texttt{\#ifdef} 噪音。

\subsubsection{\texttt{arm} 中的选择逻辑}

\begin{lstlisting}[caption={arm() selection.}, label=lst:gpuarm, style=cppstyle]
void GpuTimerPair::arm(const Device& device) {
  if (!rec_.live_ || !device.is_cuda()) return;
  if (g_gpu_trace.load(memory_order_acquire)) {
    rec_.trace_pushed_ = cupti_push_ext(static_cast<uint64_t>(rec_.slot_));
  }
  if (!g_gpu_timing.load(memory_order_acquire)) return;
  const auto stream = cuda::getCurrentCUDAStream();
  cudaEvent_t start = acquire_event(stream.device_index());
  if (!start || cudaEventRecord(start, stream.stream()) != cudaSuccess) {
    if (start) (void)cudaEventDestroy(start);
    return;
  }
  gpu_start_ = start;
  gpu_stream_ = stream.stream();
  gpu_device_ = stream.device_index();
  armed_ = true;
}
\end{lstlisting}

两个正交特征交织：

\begin{itemize}[leftmargin=1.4em]
\item \textbf{CUPTI correlation} 独立于 \texttt{g\_gpu\_timing}。如果设置了 \texttt{g\_gpu\_trace}，则 \texttt{arm} 通过 \texttt{cupti\_push\_ext} 将操作的槽作为外部相关 ID 推送到 CUPTI 的每线程堆栈上。即使事件计时关闭，也会发生这种情况，因此 \texttt{gpu\_trace} 会话关联内核而无需支付 \texttt{cudaEventRecord} 成本。
\item \textbf{cudaEvent timing} 在 \texttt{g\_gpu\_timing} 上门控。启用时，\texttt{arm} 通过 \texttt{acquire\_event} 获取池化的 \texttt{cudaEvent\_t}（从 \texttt{g\_gpu\_mutex} 下的 \texttt{g\_pool} 回收或使用 \texttt{cudaEventCreateWithFlags} 创建），将其记录在当前流上，并保存设备/流/事件。
\end{itemize}

\texttt{close} at 遵循相同的顺序：首先，如果有一个被推送（通过 \texttt{cupti\_pop\_ext} at），它会弹出 CUPTI 相关性，然后如果 \texttt{armed\_} 通过 \texttt{acquire\_event} 和 \texttt{cudaEventRecord} 记录结束事件，并在 \texttt{g\_gpu\_mutex} 下将 \texttt{LivePair\{slot,device,stream,start,end\}} 附加到 \texttt{g\_live}。分析永远不会使有效操作失败：每个 CUDA 错误路径都会破坏部分事件并返回。

\subsubsection{会话停止时的分辨率}

\texttt{gpu\_resolve\_all} 在 \texttt{profiler\_stop} 之后但在返回 Python 之前从绑定的 \texttt{\_profiler\_stop} 调用。关键步骤：

\begin{enumerate}[leftmargin=1.4em]
\item 在互斥锁下交换 \texttt{g\_live} ，以便在解析运行时可以继续调度。
\item 从 \texttt{TP\_PROFILER\_GPU\_TIMEOUT\_MS} 计算截止时间（默认 5000\,ms,）。
\item 构建 \texttt{last\_on\_stream} 映射：每个流一次等待就足够了，因为 CUDA 对流上的事件记录进行排序。
\item 对于每个流的最后一个结束事件，调用 \texttt{wait\_event} 轮询 \texttt{cudaEventQuery} 并使用 \texttt{50\,\textmu s} 休眠直到成功或截止日期。这避免了设备范围的同步。
\item 通过 \texttt{cudaEventElapsedTime} 解析每一对。成功后，写入 \texttt{events[slot].gpu\_ms = ms} 并调用 \texttt{emit} 回调。可重用事件转到本地 \texttt{reusable} 向量；失败的对将通过 \texttt{destroy\_event\_pair} 销毁，即使对于待处理的事件也会执行 \texttt{cudaEventDestroy} 。挂起的对故意不返回到池中，以便下一个会话开始干净。
\item 将可重用事件返回到一把锁下的 \texttt{g\_pool} 并恢复调用者的设备。
\end{enumerate}

该池是 \texttt{unordered\_map<int, deque<cudaEvent\_t>>}，由设备索引键控，使用 \texttt{cudaEventDefault} 创建（计时开启）。 \texttt{gpu\_drain\_pool} at 交换并销毁所有池事件和实时事件，由进程拆卸时的绑定使用。

\subsection{可选的 CUDAGuard：多 GPU 计时的正确设备}

生成的重新调度帮助程序必须在操作的目标设备上记录 GPU 事件，即使调用线程的当前设备位于其他位置也是如此。声明：

\begin{lstlisting}[caption={OptionalCUDAGuard.}, label=lst:optionalguard, style=cppstyle]
class P10_API OptionalCUDAGuard {
public:
  explicit OptionalCUDAGuard(const Device& device);
  ~OptionalCUDAGuard();
private:
  std::unique_ptr<CUDAGuard> guard_;
};
\end{lstlisting}

对于 \texttt{!USE\_CUDA} 来说，它是一个空的内联。对于 CUDA 构建，构造函数执行 \texttt{if (device.is\_cuda) guard\_ = make\_unique<CUDAGuard>(device)};析构函数通过 \texttt{CUDAGuard} 恢复。 \texttt{CUDAGuard} 本身（at）保存 \texttt{cudaGetDevice}，向目标调用 \texttt{cudaSetDevice}，并在销毁时恢复。

Codegen 在分析器处理以下位置后立即发出防护：

\begin{lstlisting}[caption={device guard emission.}, label=lst:genguard, style=cppstyle]
#ifdef USE_CUDA
    tensorplay::prof::GpuTimerPair __tp_gpu(__tp_prof_rec);
#endif
#ifdef USE_CUDA
    cuda::OptionalCUDAGuard device_guard(weight.device());
#endif
    static const OperatorHandle op_handle = Dispatcher::singleton().findHandle("add");
    DispatchKey dispatch_key = toBackendKey(...);
#ifdef USE_CUDA
    __tp_gpu.arm(weight.device());
#endif
\end{lstlisting}

放置很重要：守卫是在 \texttt{arm} 之前构造的，因此 \texttt{acquire\_event} 和 \texttt{cudaEventRecord} 在正确的设备上运行，并且 \texttt{close} 在守卫被销毁之前运行（守卫仅在函数返回之前比 \texttt{\_\_tp\_gpu} 寿命更长，但 \texttt{close} 在调度之后、销毁之前显式调用）。在 CPU 操作 \texttt{device.is\_cuda==false} 上，所以防护是无操作和 \texttt{arm} 提前返回。

\subsection{后端堆栈：ITT、NVTX 和 CUPTI}

\begin{table}[H]
\centering
\caption{探查器后端。编译时启用；在运行时通过原子选择。}
\label{tab:profbackends2}
\small
\begin{tabular}{llll}
\toprule
后端 & File & 机制 & 运行时库 \\
\midrule
CPU & \texttt{Profiler.cpp} & \texttt{chrono::steady\_clock}，矢量漏极 & -- \\
GPU 事件 & \texttt{ProfilerGpu.cpp} & \texttt{ElapsedTime} & \texttt{cudart} \\
中国电力工业联合会 & \texttt{ProfilerCupti.cpp} & \texttt{cuptiActivityEnable}，缓冲区回调 & \texttt{.11} 通过 \texttt{dlopen} \\
Stub & \texttt{ProfilerCuptiStub.cpp} & 无操作符号 & --（CPU/HIP 构建） \\
NVTX & \texttt{ProfilerNvtx.cpp} & \texttt{Pop} & \texttt{libnvtx.so} 通过 \texttt{dlopen} \\
ITT & \texttt{ProfilerItt.cpp} & \texttt{end} & \texttt{libittnotify.so} 通过 \texttt{dlopen} \\
\bottomrule
\end{tabular}
\end{table}

\subsubsection{分析器NVTX}

NVTX 桥是 dlopen'd，未链接。 \texttt{ensure\_loaded} 位于探头 \texttt{libnvtx.so}、\texttt{libnvtx\_dynamic.so}、\texttt{libNVToolsExt.so} 和 \texttt{libroctx64.so}（ROCm 双胞胎，导出 \texttt{roctxRangePushA} 等），通过 \texttt{dlopen} 和 \texttt{dlsym} 实现 \texttt{nvtxRangePushA}、\texttt{nvtxRangePop}、\texttt{nvtxMarkA}、\texttt{nvtxRangeStartA}、\texttt{nvtxRangeEnd}。如果没有加载，\texttt{g\_available=false} 和每个调用都会降级。存在两个条目组：

\begin{itemize}[leftmargin=1.4em]
\item \texttt{tensorplay.cuda.nvtx} \texttt{range\_end} 的原始直通。当不可用时，它们会抛出 \texttt{RuntimeError} 并显示“未安装 NVTX 功能”，从而保留历史合约。
\item 内部无抛出钩子 \texttt{end} 检查 \texttt{emitting = g\_emit\_nvtx.load(relaxed) \&\& usable} 然后调用 \texttt{ nv\_range\_pop}。 \texttt{OpRecord::begin} 仅当 \texttt{nvtx\_on} 为 true 时才调用它们（这还需要 \texttt{kind==kOp}），因此用户注释不会发出 NVTX，除非通过绑定的 \texttt{\_profiler\_emit\_nvtx} 明确请求。
\end{itemize}

\texttt{g\_emit\_nvtx} 通过 \texttt{emit\_nvtx} 上下文管理器从 Python 切换。

\subsubsection{分析器Itt}

ITT 桥接器遵循与 NVTX 相同的 dlopen 结构，但针对 Intel VTune/Advisor 系列。 \texttt{ensure\_loaded} 在 dlopens \texttt{libittnotify.so} 处，然后 \texttt{libittnotify64.so}、dlsyms \texttt{\_\_itt\_domain\_create}、\texttt{\_\_itt\_string\_handle\_create}、\texttt{\_\_itt\_task\_begin}、\texttt{\_\_itt\_task\_end} 创建域 \texttt{"tensorplay"}。内部钩子 \texttt{end} 检查 \texttt{g\_emit\_itt.load(relaxed)} 然后调用 \texttt{end} ，它在内部执行 \texttt{fn\_string\_handle\_create} + \texttt{fn\_task\_begin} 。可用性为 \texttt{itt\_available}。

两个桥均不保留硬链接依赖性：仅 CPU 轮导入并运行，所有探查器调用都变为无操作。

\subsubsection{库蒂存根}

对于 HIP 版本（以及任何重新导出的 \texttt{ProfilerCuptiStub.cpp}），该文件为 \texttt{g\_gpu\_trace}、\texttt{cupti\_available}、\texttt{cupti\_last\_error}、\texttt{cupti\_start}、\texttt{cupti\_stop\_and\_collect}、\texttt{cupti\_push\_ext}、\texttt{cupti\_pop\_ext}、\texttt{cupti\_version} 提供无操作定义。这使得扩展可以在 \texttt{RTLD\_NOW} 下导入，而无需嵌入 \texttt{\#ifdef} 引用 CUPTI 符号（\texttt{ProfilerGpu.cpp}、\texttt{init.cpp}）的约 15 个调用站点。 CPU 版本在 \texttt{!USE\_CUDA} 时提供相同的后备。

\subsection{CUPTI 详细信息：订阅、缓冲区和correlationId}

\texttt{ProfilerCupti.cpp} 对于 USE\_CUDA 是 596 行，对于 CPU 是 24 行。设计约束是： \texttt{libcupti} 绝不是硬依赖项；所有入口点都是从 dlopened 句柄（\texttt{108,424--466}）进行 dlsym 处理，并且没有工具包的计算机上的会话会降级为 \texttt{cupti\_available==false}。

\subsubsection{入口点和句柄}

\texttt{CuptiFns} 保存 \texttt{GetTimestamp}、\texttt{RegisterCallbacks}、\texttt{Enable}、\texttt{Disable}、\texttt{FlushAll}、\texttt{GetNextRecord}、\texttt{GetNumLostRecords}、\texttt{PushExternalCorrelationId}、\texttt{PopExternalCorrelationId}、\texttt{GetCallbackName}、\texttt{GetVersion}。 \texttt{g\_cupti\_lib} 在进程生命周期内保留； \texttt{g\_fns} 在第一次成功 \texttt{cupti\_start} 之后设置一次，因此 \texttt{pop} 可以在不锁定的情况下读取它（注释）。

\texttt{cupti\_available} 在未启用的情况下进行探测：它首先检查 \texttt{g\_fns!=nullptr}，然后尝试 \texttt{dlopen(..., RTLD\_NOLOAD)} 查找 \texttt{.so}。 \texttt{cupti\_version} at dlsyms \texttt{cuptiGetVersion} 不启用任何类型，用于标记跟踪导出模式元数据。

\subsubsection{订阅：种类和回调}

\texttt{cupti\_start} at（从 \texttt{profiler\_start} 之前调用）执行以下操作：

\begin{enumerate}[leftmargin=1.4em]
\item 通过宏 \texttt{TP\_CUPTI\_DLSYM} 从候选 \texttt{libcupti.so.13,.12,.11,.1,.so} 和 dlsyms 强制符号中延迟 dlopens 库；可选符号 \texttt{GetNumLostRecords}、\texttt{GetCallbackName}、\texttt{GetVersion} 进行空检查。
\item 校准时基：\texttt{g\_fns->GetTimestamp(\&cupti\_now)} 然后 \texttt{g\_time\_offset\_ns = steady\_now\_ns - cupti\_now}。每个活动记录的 \texttt{end} 稍后都会通过该值进行偏移，将 CUPTI 的历元映射到操作时间线的 \texttt{CLOCK\_MONOTONIC} 上。
\item 通过 \texttt{RegisterCallbacks} 注册缓冲区回调 \texttt{buffer\_requested} / \texttt{buffer\_completed}。
\item 启用五种：
\begin{lstlisting}[caption={enabled kinds.}, label=lst:cuptiKinds, style=cppstyle]
static const CUpti_ActivityKind kinds[] = {
  CUPTI_ACTIVITY_KIND_MEMCPY,
  CUPTI_ACTIVITY_KIND_MEMSET,
  CUPTI_ACTIVITY_KIND_RUNTIME,
  CUPTI_ACTIVITY_KIND_CONCURRENT_KERNEL,
  CUPTI_ACTIVITY_KIND_EXTERNAL_CORRELATION,
};
for (auto k : kinds) rc = g_fns->Enable(k);
\end{lstlisting}
\end{enumerate}

\texttt{disable\_kinds\_locked} at 禁用 \texttt{cupti\_stop\_and\_collect} 中的相同设置。

缓冲区处理在 (\texttt{kBufSize = 1<<20}) 处使用 1\,MiB 池缓冲区：

\begin{itemize}[leftmargin=1.4em]
\item \texttt{buffer\_requested} 从 \texttt{g\_cupti\_mutex} 下的 \texttt{g\_free\_bufs} 弹出，或者执行 \texttt{malloc(kBufSize)}。分配失败时，它会返回 \texttt{0}，因此 CUPTI 会删除记录，但分析会继续。
\item \texttt{buffer\_completed} 由 CUPTI 在其完成线程上调用；它锁定 \texttt{g\_cupti\_mutex}，调用 \texttt{parse\_buffer\_locked}，然后将缓冲区推送到 \texttt{g\_free\_bufs} 进行回收。
\item \texttt{parse\_buffer\_locked} 通过 \texttt{GetNextRecord} 进行迭代。对于 \texttt{CONCURRENT\_KERNEL}，它选择修订版别名 \texttt{ActivityKernel}（CUDA 12.8+ 上的 Kernel9、Kernel8 或 Kernel4），并通过 \texttt{intern\_name\_locked} 推送名称为 \texttt{GpuActivity\{'k',\dots\}} 的 \texttt{GpuActivity\{'k',\dots\}}。 Memcpy/memset/runtime/driver/external-correlation 每个都映射到 \texttt{'d'} 加上注释中的前缀字段提取（修订仅附加尾随字段）。丢失的记录通过 \texttt{GetNumLostRecords} 显示。
\end{itemize}

\subsubsection{correlationId 和外部关联}

关键连接是三向的：op $\rightarrow$ 运行时 API $\rightarrow$ 内核。 \texttt{GpuTimerPair::arm} 将操作的槽位作为 \texttt{uint64\_t} 外部 id 推送； CUPTI 发出 \texttt{EXTERNAL\_CORRELATION} 记录映射 \texttt{correlationId -> externalId} (, kind \texttt{CUSTOM0}); kernel/memcpy 活动记录带有相同的 \texttt{correlationId}。在快照时间 \texttt{cupti\_stop\_and\_collect} 通过 \texttt{FlushAll(0)} 刷新，休眠 5ms 以吸收完成线程，然后在以下位置将 \texttt{g\_acts} 与 \texttt{g\_corr2ext} 连接：

\begin{lstlisting}[caption={external-id join.}, label=lst:correlationJoin, style=cppstyle]
for (auto& a : acts) {
  auto it = g_corr2ext->find(a.correlation);
  if (it != g_corr2ext->end()) a.external_id = it->second;
}
\end{lstlisting}

然后每个操作的绑定都会折叠：

\begin{lstlisting}[caption={per-op kernel aggregation.}, label=lst:foldKernels, style=cppstyle]
for (auto& a : gpu_acts) {
  if (a.kind=='r' || a.kind=='d') continue;
  if (a.external_id==GpuActivity::kNoExt || a.external_id>=events.size()) continue;
  auto& op = events[a.external_id];
  double ms = double(a.end_ns - a.start_ns)/1e6;
  op.gpu_ms = (op.gpu_ms < 0.f ? 0.f : op.gpu_ms) + float(ms);
  op.kernel_count += 1;
}
\end{lstlisting}

\texttt{pop\_ext} 是 \texttt{CUPTI\_EXTERNAL\_CORRELATION\_KIND\_CUSTOM0} 上的每线程堆栈操作；复合内部操作正确嵌套，因为 \texttt{GpuTimerPair} 将每个调度括起来。在 \texttt{gpu\_trace} 会话之外，唯一的每操作成本是 \texttt{GpuTimerPair::arm} （注释）中的 \texttt{g\_gpu\_trace} 获取加载，与 \texttt{g\_gpu\_timing} 相同。

\subsection{集成一：代码生成}

每个分派的执行都恰好经过一个 \texttt{detail::redispatch\_*} 漏斗。生成器工具会漏斗，以便分析器在 autograd 和 autocast 下面看到每个操作一次。在 (\texttt{\_emit\_redispatch})：

\begin{lstlisting}[caption={redispatch instrumentation.}, label=lst:genredispatch, style=cppstyle]
namespace detail {
TENSORPLAY_API Tensor redispatch_add(const Tensor& self, const Tensor& other) {
    tensorplay::prof::OpRecord __tp_prof_rec("add");
    if (tensorplay::prof::g_capture_shapes.load(std::memory_order_acquire)) {
        std::vector<std::vector<int64_t>> __tp_prof_shapes;
        std::vector<int32_t> __tp_prof_dtypes;
        __tp_prof_shapes.push_back(static_cast<std::vector<int64_t>>(self.shape()));
        __tp_prof_dtypes.push_back(static_cast<int32_t>(self.dtype()));
        __tp_prof_shapes.push_back(static_cast<std::vector<int64_t>>(other.shape()));
        __tp_prof_dtypes.push_back(static_cast<int32_t>(other.dtype()));
        __tp_prof_rec.set_io_meta(std::move(__tp_prof_shapes), std::move(__tp_prof_dtypes));
    }
#ifdef USE_CUDA
    tensorplay::prof::GpuTimerPair __tp_gpu(__tp_prof_rec);
#endif
#ifdef USE_CUDA
    cuda::OptionalCUDAGuard device_guard(self.device());
#endif
    static const OperatorHandle op_handle = Dispatcher::singleton().findHandle("add");
    DispatchKey dispatch_key = toBackendKey(dispatchKeyForTensorArgs(self, other));
#ifdef USE_CUDA
    __tp_gpu.arm(self.device());
#endif
    auto&& __tp_result = DispatchStub<Tensor, const Tensor&, const Tensor&>::call(op_handle, dispatch_key, self, other);
#ifdef USE_CUDA
    __tp_gpu.close();
#endif
    if (tensorplay::prof::g_capture_shapes.load(std::memory_order_acquire)) {
        __tp_prof_rec.set_output_bytes(static_cast<int64_t>(__tp_result.numel()) * static_cast<int64_t>(__tp_result.itemsize()));
    }
    return std::move(__tp_result);
}
} // namespace detail
\end{lstlisting}

细节：

\begin{itemize}[leftmargin=1.4em]
\item \texttt{\_\_tp\_prof\_rec} 首先在形状分支之前构造，因此即使未捕获形状也会发生获取加载。形状分支迭代类似张量的参数（\texttt{is\_tensor\_like}，处理\texttt{is\_list}，\texttt{is\_opt}，\texttt{is\_mutable\_ref}），并仅在设置\texttt{g\_capture\_shapes}时复制形状/数据类型。分支内的第二个负载意味着完全不活动的路径恰好消耗一个负载。
\item 对于可变操作，保留 \texttt{DispatchStub::call} 之后的版本冲突；对于 \texttt{void} 返回形状/输出字节尾部被省略。
\item \texttt{set\_output\_bytes} 在相同的 \texttt{g\_capture\_shapes} 负载上门控，并且仅适用于 \texttt{Tensor} 返回操作。该值提供内存快照视图，与分配器级 \texttt{MemEvent} 记帐不同。
\item 生成器标头包括 \texttt{Profiler.h} 处的 \texttt{Profiler.h} 和 \texttt{USE\_CUDA} 下的 \texttt{CUDARuntime.h}，因此发出的 TU 可在 CPU 和 CUDA 版本上进行编译。
\end{itemize}

处理 vmap、autocast 和 autograd 调度键后，委托到 \texttt{detail::redispatch\_add} 的公共 \texttt{Tensor::add} 条目；因此，相同的 \texttt{OpRecord} 涵盖了急切执行和图形捕获执行，并且仍然通过调度程序路由的融合 \texttt{stax} 操作也报告字节。

\subsection{集成二：Autograd 引擎}

引擎发出两个跨度：

\subsubsection{每个节点的向后事件}

\begin{lstlisting}[caption={per-node backward span.}, label=lst:backwardspan, style=cppstyle]
std::optional<tensorplay::prof::OpRecord> __tp_node_rec;
if (tensorplay::prof::g_active.load(std::memory_order_acquire)) {
  __tp_node_rec.emplace(tensorplay::prof::intern_name("backward::" + func->name()));
}
\end{lstlisting}

它位于 \texttt{Engine::evaluate\_function} 内部。仅当活动时，跳跃成本是一个原子负载加上一个虚拟 \texttt{func->name} demangle 和 \texttt{intern\_name} dedup（每个不同节点名称一个 arena 副本）。 \texttt{optional<OpRecord>} 确保析构函数甚至在异常路径上运行； \texttt{AnomalyMode} 处理保留了失败回溯。节点事件按墙时间与父 \texttt{\_\_backward\_\_} 跨度重叠，在 chrome 跟踪中作为子节点可见。

\subsubsection{后向相位跨度}

\begin{lstlisting}[caption={backward-phase annotation.}, label=lst:backwardsession, style=cppstyle]
tensorplay::prof::OpRecord __tp_backward_span("__backward__", tensorplay::prof::EventKind::kBackward);
\end{lstlisting}

它位于 \texttt{Engine::execute} 中。它包含整个 \texttt{execute} 调用，包括工作线程可能同时评估设备节点的 \texttt{while(!graph\_task.is\_completed)} 消耗（设备 \texttt{ReadyQueue} 按 \texttt{sequence\_nr} 确定优先级）。在 \texttt{create\_graph} 下创建的二阶图连接，因为 \texttt{GradModeGuard} 在 \texttt{evaluate\_function} 期间处于活动状态。

\subsection{集成三：Python 绑定 Drain 和 Trace Export}

连接 C++ 和 Python：

\begin{table}[H]
\centering
\caption{Python 绑定探查器符号。}
\label{tab:pyprof}
\small
\begin{tabular}{lll}
\toprule
象征 & 线路 & Role \\
\midrule
\texttt{\_profiler\_start} &  & 存储 \texttt{g\_gpu\_trace}，调用 \texttt{cupti\_start}，然后调用 \texttt{profiler\_start*} \\
\texttt{\_profiler\_stop} &  & \texttt{profiler\_stop}、\texttt{gpu\_resolve\_all}、\texttt{cupti\_stop\_and\_collect}、\texttt{mem\_take}、元组转换 \\
\texttt{\_profiler\_is\_active} &  & 用于 Python 端跨度发射的单次获取加载 \\
\texttt{end} &  & 捕获 Python 站点然后 \texttt{end} \\
\texttt{itt} &  & 存储到 \texttt{g\_emit\_itt} \\
\texttt{\_profiler\_cupti\_version} &  & 标记跟踪模式元数据 \\
\bottomrule
\end{tabular}
\end{table}

\texttt{\_profiler\_start} 首先存储 \texttt{g\_gpu\_timing}（第 501 行），因此稍后的 CPU 配置文件不会为 CPU 重新分派准备事件，然后如果 \texttt{gpu\_trace} 在第一个操作之前准备 CUPTI（第 509--519 行），并在不可用时通过 \texttt{PyErr\_WarnEx} 发出警告。然后分支到 \texttt{ profiler\_start} （第 521--527 行）。

\texttt{\_profiler\_stop} at 围绕停止工作释放 GIL（第 540 行），通过有界等待解析 GPU 对（第 545 行），原子交换 \texttt{g\_gpu\_trace}（第 550 行），收集 CUPTI 活动，按操作折叠它们（第 557--569 行），获取 mem 事件（第 571--572 行），重新获取 GIL，然后将每个 \texttt{Event} 转换为第 574--610 行的 13 元组 \texttt{(name,kind,start\_ns,end\_ns,tid,shapes,dtypes,site,gpu\_ms,out\_bytes,stack,kernel\_count,flops)}。通过 \texttt{estimate\_flops} 处的热路径中的每个事件计算一次 FLOP（矩阵乘法算作两个操作，转换以步幅 1 / 填充 0 近似）。

(\texttt{class profile}) 处的高级 Python 界面编排会话：

\begin{lstlisting}[caption={session start/stop.}, label=lst:profilerpy, style=pystyle]
def _start_session(self):
    _C._profiler_start(self.record_shapes, self.with_stack, self.gpu_timing, self.gpu_trace, self.profile_memory)
    self._recording = True
    _library._note_profiling_start()
def _stop_session(self):
    raw_ops, raw_gpu, raw_mem = _C._profiler_stop()
    events = self._ensure_events()
    events.extend(raw_ops)
    self.gpu_activities.extend(raw_gpu)
    self.mem_events.extend(raw_mem)
    self._recording = False
\end{lstlisting}

\texttt{profile} 支持 \texttt{schedule} 回调（\texttt{ NONE} 位于 \texttt{\_schedule.py}），并在 \texttt{with\_modules=True}（发出 \texttt{\_profiler\_user\_begin("nn.Module:...")}）时通过修补 \texttt{Module.\_\_call\_\_} 进行模块跟踪。委托处的 Chrome 跟踪导出，通过 \texttt{correlationId} 合并 CPU 和 GPU 事件并应用 CUPTI 偏移量。

\subsection{用数字进行开销分析}

\begin{table}[H]
\centering
\caption{探查器开销预算。在安静主机、中档 x86、单线程、\texttt{bench\_profiler\_overhead.py} 上测量。}
\label{tab:profoverhead}
\small
\begin{tabular}{llll}
\toprule
Mode & 每个操作运行什么 & Cost & 笔记 \\
\midrule
残疾人 & \texttt{g\_active.load(acquire)} + 未采用的分支 & $\approx$ 2 个周期（$\sim$0.6\,ns，3.5\,GHz） & 一次负载是唯一的热路径成本 \\
启用 CPU（无形状） & \texttt{chrono::steady\_clock} + 向量推 & $\approx$ 30\,ns & \texttt{now\_ns} at + \texttt{g\_mutex} + \texttt{push\_back} \\
CPU + 形状 & 上面+形状/dtype矢量副本 & $\approx$ 80--120\,ns & 第二个 \texttt{g\_capture\_shapes} 加载于 + \texttt{set\_io\_meta} \\
CPU + 带\_stack & 以上+现场实习生 & $\approx$ 150--300\,ns & \texttt{intern\_site} 在 + \texttt{t\_pending\_site} 采用 \\
GPU 事件（cudaEvent） & 每个操作两个 \texttt{cudaEventRecord} & $\approx$ 5\,\textmu s & 因此通过 \texttt{gpu\_timing} 选择加入；不适用于细粒度内核 \\
CUPTI 跟踪 & CUPTI 缓冲区填充 + 完成线程上的解析 & $\approx$ 1--3\,\textmu 摊销 & 1\,MiB 缓冲区；解析于 \\
NVTX/ITT & 一范围推/弹出 & $\approx$ 40--60\,ns 发射时 & 由 \texttt{itt} 松弛负载保护 \\
\bottomrule
\end{tabular}
\end{table}

\subsubsection{为什么禁用是2个周期}

\texttt{OpRecord::begin} 是静态原子的一个 \texttt{memory\_order\_acquire} 负载，内联到 x86 上的 \texttt{mov} + \texttt{cmp} + 预测未采用分支。除了负载之外没有线程本地，没有调用。形状/站点第二个负载位于 \texttt{g\_active} 真分支内部，因此在禁用时它们会添加零周期。这与每个操作周围已经发出的 \texttt{GradMode::is\_enabled} 和 \texttt{InferenceMode::is\_enabled} 防护相匹配。

\subsubsection{为什么CPU是30\,ns}

\texttt{now\_ns} 是 \texttt{chrono::steady\_clock::now.time\_since\_epoch.count}，在 Linux 上是 \texttt{clock\_gettime(CLOCK\_MONOTONIC)} vdso 调用 ($\sim$15\,ns)。剩余的 $\sim$15\,ns 是 \texttt{g\_mutex} 无竞争锁加上 \texttt{g\_events->push\_back} 摊销分配。在 256$\times$512（单线程固定，\texttt{--threads 1 --framework tp}）上对 \texttt{add} 进行 200 次迭代时测得的 p50 增加了 28--35\,ns 延迟。

\subsubsection{为什么 GPU 事件是 5\,\textmu s}

\texttt{cudaEventRecord} 是一个驱动程序调用，它将标记放入流中；每个操作对花费两条记录加上 \texttt{acquire\_event} 池查找。在引用的 sm\_89 远程盒上，测量的往返时间为 4.2--6.1\,\textmu s。这就是默认情况下不启用 GPU 计时的原因：它支配短于 $\sim$10\,\textmu s 的内核。对于 sub-10\,\textmu 内核，\ac{cupti} 的 \texttt{CONCURRENT\_KERNEL} 活动更可取，因为它在 GPU 上添加时间戳并异步报告，无需每次操作驱动程序调用。

使用 \texttt{--threads 1} 时，CPU 开销列在 256$\times$512 matmul 上应为 8--14\%（基础 $\sim$12\,\textmu s，分析 $\sim$13.2\,\textmu s），在内核时间占主导地位的较大形状上缩小到 $<$2\%。

\subsection{互动之旅}

\begin{lstlisting}[caption={Profiling a few ops from Python (no GPU required).}, label=lst:smoketest, style=pystyle]
python3 - << 'PY'
import tensorplay as tp
from tensorplay.profiler import profile, ProfilerActivity
with profile(record_shapes=True, with_stack=True) as prof:
    a = tp.ones([8, 16])
    b = tp.ones([8, 16])
    c = tp.add(a, b)
    d = c.relu()
print(prof.key_averages().table(sort_by="cpu_time_total", row_limit=10))
# GPU events require a CUDA build; degraded gracefully on CPU:
with profile(activities=[ProfilerActivity.CPU]) as prof:
    x = tp.randn([32, 32])
    y = x @ x
print("events:", len(prof.events))
PY

# chrome trace export
python3 - << 'PY'
import tensorplay as tp
from tensorplay.profiler import profile
with profile(record_shapes=True) as prof:
    for _ in range(5):
        tp.add(tp.ones([4,4]), tp.ones([4,4]))
prof.export_chrome_trace("/tmp/tp_trace.json")
print(open("/tmp/tp_trace.json").read()[:800])
PY
cat /tmp/tp_trace.json | python3 -m json.tool | head -n 40
\end{lstlisting}

\subsection{不变量和边缘情况}

\begin{itemize}[leftmargin=1.2em]
\item \textbf{A1: Exactly one OpRecord per dispatched op.} 重新调度漏斗是 autograd/autocast/vmap 下面的单一检测站点。 Public \texttt{Tensor::add} 从不直接检测；它委托给 \texttt{detail::redispatch\_add}。不通过 codegen 调用 \texttt{Dispatcher::singleton.findHandle} 的自定义运算符必须发出自己的 \texttt{OpRecord} 或不可见（\texttt{library.py} 处的自定义操作包装器在 \texttt{\_note\_profiling\_start} 处于活动状态时执行此操作，\texttt{profiler.py}）。
\item \textbf{A2: Slots stay stable until next start.} \texttt{g\_events} 仅在会话开始时清除，并在会话停止时移出；名称 arena 也在那里被清除，但在每个插槽移动之后不会被清除，保留借用的指针直到绑定复制它们。
\item \textbf{A3: GPU pair stranded on exception is destroyed, not leaked.} 如果 \texttt{DispatchStub::call} 在 \texttt{close} 之前抛出，则 \texttt{GpuTimerPair} 析构函数调用 \texttt{close}，后者仍然执行 \texttt{cupti\_pop\_ext} 并在 \texttt{!rec\_.live\_} 时销毁任何获取的开始事件。
\item \textbf{A4: Stop waits for inflight.} 会话停止在交换 \texttt{g\_events} 之前等待 \texttt{g\_inflight\_cv}，因此 \texttt{gpu\_resolve\_all} 永远不会与仍在展开的 \texttt{\textasciitilde OpRecord} 竞争，后者会在快照之后写入 \texttt{end\_ns}。
\item \textbf{A5: CUPTI FlushAll runs without the CUPTI mutex.} 互斥体在 \texttt{FlushAll(0)} 之前释放，因为 \texttt{buffer\_completed} 使用相同的互斥体来解析和回收缓冲区；持有它会陷入僵局。
\item \textbf{A6: No hard CUPTI/NVTX/ITT link.} 所有三个桥 dlopen,。缺少的库会降级为 \texttt{cupti\_available==false} / \texttt{nvtx\_available==false} / \texttt{itt\_available==false} 并带有警告，而不是链接错误。存根为 HIP 构建提供相同的接口。
\end{itemize}

